\scanpage{83}
- \mspage{IV-1014} -
module de $P[T_1,\ldots,T_n]/((Q\cap P)[T_1,\ldots,T_n])=(P/(Q\cap P))[T_1,\ldots,T_n]$. En tant
que $(A/\underline L)$-module, il est donc isomorphe à un sous-module d'une somme
directe de $(A/\underline L)$-modules isomorphes à $P/(Q\cap P)=(P+Q)/Q$. Mais $(P+Q)/Q$
est un $(A/\underline L)$-module isomorphe à un sous-module de $P[T_1,\ldots,T_n]/Q$, qui
n'est autre que $\operatorname{gr}_{\underline L}^\bullet(M)$ en vertu du diagramme (19.7.1.1). On a donc
\[
\text{(19.7.1.4)}\qquad
\operatorname{Ass}_{A/\underline L}(Z)\subset
 \operatorname{Ass}_{A/\underline L}(\operatorname{gr}_{\underline L}^\bullet(M)).
\]
Mais chacun des $\operatorname{gr}_{\underline L}^m(M)$ est par hypothèse un $(A/\underline L)$-module \uline{plat} de type
fini, donc projectif puisque $A/\underline L$ est noethérien, et on a par suite
$\operatorname{Ass}_{A/\underline L}(\operatorname{gr}_{\underline L}^\bullet(M))
 \subset\operatorname{Ass}_{A/\underline L}(A/\underline L)$.
Pour voir que $Z=0$, il suffit par suite (Bourbaki, \uline{Alg.comm.}, chap.IV, \S~1,
n$^\circ$1, cor.1 de la prop.2\inserted{ et n$^\circ$3, prop.7}) de voir que pour
tout \deleted{\illegible{3} $\unsure{\underline p}$ \illegible{2}}
$\underline q\in\operatorname{Ass}_{A/\underline L}(A/\underline L)$, on
\deleted{\illegible{1}} a $Z_{\underline q}=0$.

Mais les idéaux de $\operatorname{Ass}_{A/\underline L}(A/\underline L)$ sont les idéaux de la forme
$\underline p/(\underline L/\underline K)$ où
$\underline p\in\operatorname{Ass}_{A/\underline K}(A/\underline L)$, et il revient donc au même de voir que
$Z_{\underline p}=0$ pour tout
$\underline p\in\operatorname{Ass}_{A/\underline K}(A/\underline L)$.
Or, si $\underline p=\underline r/\underline K$, où $\underline r$ est un idéal premier de $A$, on a
$(\operatorname{gr}_{\underline K}^\bullet(M))_{\underline p}
 =\operatorname{gr}_{\underline K_{\underline r}}^\bullet(M_{\underline r})$
et $(A/\underline K)_{\underline p}=A_{\underline r}/\underline K_{\underline r}$,
et les images des $f_i$ dans $A_{\underline r}$ forment
une suite $\operatorname{gr}_{\underline K_{\underline r}}^\bullet(A_{\underline r})$-régulière
(\uline{0},15.1.14); appliquant à $A_{\underline r},\underline K_{\underline r}$ et
$M_{\underline r}$ l'implication a)$\Rightarrow$b) de l'énoncé, on voit que
\deleted{\illegible{5}}
l'hypothèse d) entraîne que $\psi_{M_{\underline r}}$ est bijectif, donc aussi
$\psi'_{R_{\underline r}}$ en vertu
de (19.6.3); en d'autres termes, $Q_{\underline r}=Q'_{\underline r}[T_1,\ldots,T_n]$,
où on a posé
$Q'=\operatorname{gr}_{\underline K}^\bullet(R,N)\otimes_{A/\underline K}(A/\underline L)$,
qui est ici un sous-module de $P$; en particulier
$Q'_{\underline r}=Q_{\underline r}\cap P_{\underline r}$, donc finalement on a bien
$Z_{\underline r}=Z_{\underline p}=0$, ce qui achève de
démontrer (19.7.1.3).

(19.7.1.5) En vertu de (19.7.1.3), pour prouver que $\psi'_R$ est surjectif,
il suffit de montrer que tout élément homogène de degré $m$ de $Q\cap P$ est
l'image d'un élément de même degré de
$\operatorname{gr}_{\underline K}^\bullet(R,N)\otimes_{A/\underline K}(A/\underline L)$.
Or, les éléments
de degré $m$ de $P$ s'identifient\deleted{\unsure{canoniquement}}
(par $\psi_N$) aux éléments
\deleted{\unsure{homogènes}} de $\operatorname{gr}_{\underline L}^m(N)$
qui appartiennent à
$(\underline K^mN+\underline L^{m+1}N)/\underline L^{m+1}N$; ceux
qui sont images d'éléments \deleted{\unsure{homogènes}} de
$\operatorname{gr}_{\underline K}^m(R,N)\otimes_{A/\underline K}(A/\underline L)$ s'identi
\mspage{82}
